\begin{explanationbox}
	\textbf{\textcolor{CExplanation}{一句话直觉}}

	两个独立平方和之积，总不小于它们交叉项和的平方；其本质是向量的长度不小于点积的绝对值.

	\medskip
	\textbf{\textcolor{CExplanation}{核心拆解}}
	\begin{description}[style=nextline, leftmargin=2.8em, labelsep=0.5em, itemsep=0.45em]
		\item[\textbf{要点一（平方结构）：}] 左边是 $(a^2+b^2)(c^2+d^2)$，右边是 $(ac+bd)^2$，两者通过完全平方差相联系.
		\item[\textbf{要点二（等号条件）：}] 等号成立要求 $ad=bc$，即对应分量成比例.
	\end{description}

	\medskip
	\textbf{\textcolor{CExplanation}{几何本质（向量长度点积）}}

	设向量 $\boldsymbol{u}=(a,b),\ \boldsymbol{v}=(c,d)$，则不等式化为 $|\boldsymbol{u}|^2|\boldsymbol{v}|^2\ge (\boldsymbol{u}\cdot\boldsymbol{v})^2$，等价于 $|\boldsymbol{u}||\boldsymbol{v}|\ge |\boldsymbol{u}\cdot\boldsymbol{v}|$。这正是 $|\cos\theta|\le 1$，其中 $\theta$ 为两向量夹角.

	\medskip
	\textbf{\textcolor{CExplanation}{代数意义（判别式非负）}}

	构造关于 $x$ 的二次函数 $f(x)=(a^2+b^2)x^2-2(ac+bd)x+(c^2+d^2)$，其值恒非负，故判别式 $\Delta=4[(ac+bd)^2-(a^2+b^2)(c^2+d^2)]\le0$，即得不等式.

	\medskip
	\textbf{\textcolor{CExplanation}{考点价值（常见题型）}}
	\begin{itemize}[leftmargin=*, itemsep=0.5em, topsep=0.4em]
		\item \textbf{考法一（求最值）：} 直接套用柯西不等式求二元表达式的最大值或最小值.
		\item \textbf{考法二（证明不等式）：} 将目标式放缩为柯西不等式的形式进行证明.
	\end{itemize}

	\medskip
	\textbf{\textcolor{CExplanation}{顿悟点}}

	柯西不等式将四个数的乘积转化为和的平方，结构对称。使用时需要识别平方和与内积的形式，并验证等号条件是否可达到.

	\medskip
	\textbf{\textcolor{CExplanation}{使用场景}}
	\begin{itemize}[leftmargin=*, itemsep=0.5em, topsep=0.4em]
		\item \textbf{场景一（求最值）：} 已知 $a^2+b^2$ 和 $c^2+d^2$ 或直接给出平方和时，求 $ac+bd$ 的最值.
		\item \textbf{场景二（不等证明）：} 在证明中构造出 $(a^2+b^2)(c^2+d^2)-(ac+bd)^2$ 的差，并利用其为完全平方.
	\end{itemize}
\end{explanationbox}
